\documentclass{article}
\usepackage{amsmath}
\usepackage{amssymb}
\usepackage{amsthm}

\input{../../preamble/layout-letter.tex}

\newtheorem{theorem}{Theorem}
\theoremstyle{definition}
\newtheorem{definition}{Definition}

\begin{document}

\begin{enumerate}
  \item Let $A, B, C$, and $D$ be subsets of a universal set $\mathcal{U}$. 
  Prove the following propositions.\\
  \textit{This proves that the three subsets are a partition of $(A \cup B)$. Partitions will be very important to understanding equivalence relations.}
  \begin{enumerate}
  \item $(A \cup B)=(A \cap B) \cup(A \backslash B) \cup(B \backslash A)$.
  \begin{proof}
    By definition, the equality $(A \cup B)=(A \cap B) \cup(A \backslash B) \cup(B \backslash A)$ 
    requires showing both directions of subset relations, 
    $(A \cup B) \subseteq (A \cap B) \cup(A \backslash B) \cup(B \backslash A)$ and 
    $(A \cup B) \supseteq (A \cap B) \cup(A \backslash B) \cup(B \backslash A)$.\vspace*{10px}\\
    $(\subseteq)$ Let $x \in A \cup B$ be arbitrary. Then $x \in A$ or $x \in B$ 
    by definition of union. Consider the following cases.
    \begin{enumerate}
    \item If $x \in A$ and $x \in B$, then $x \in A \cap B$.
    \item If $x \in A$ and $x \notin B$, then $x \in A \backslash B$.
    \item If $x \notin A$ and $x \in B$, then $x \in B \backslash A$.
    \end{enumerate}
    In all cases, $x \in (A \cap B) \cup(A \backslash B) \cup(B \backslash A)$. 
    Hence, $(A \cup B) \subseteq (A \cap B) \cup(A \backslash B) \cup(B \backslash A)$.\vspace*{10px}\\
    $(\supseteq)$ Let $x \in (A \cap B) \cup(A \backslash B) \cup(B \backslash A)$ be arbitrary. 
    Then $x \in A \cap B$ or $x \in A \backslash B$ or $x \in B \backslash A$. Consider the following cases.
    \begin{enumerate}
    \item If $x \in A \cap B$, then $x \in A$ and $x \in B$, so $x \in A \cup B$.
    \item If $x \in A \backslash B$, then $x \in A$, so $x \in A \cup B$.
    \item If $x \in B \backslash A$, then $x \in B$, so $x \in A \cup B$.
    \end{enumerate}
    Therefore, $(A \cap B) \cup(A \backslash B) \cup(B \backslash A) \subseteq A \cup B$.\vspace*{10px}\\
    Therefore, by $(\subseteq)$ and $(\supseteq)$, 
    $(A \cup B)=(A \cap B) \cup(A \backslash B) \cup(B \backslash A)$.
  \end{proof}

  \item $(A \cap B)$, $(A \backslash B)$, and $(B \backslash A)$ are all pairwise disjoint.
  \begin{proof}
    Show that any two of these sets have empty intersection.
    \begin{enumerate}
    \item To claim $(A \cap B) \cap (A \backslash B) = \emptyset$, 
    assume there exists $x \in (A \cap B) \cap (A \backslash B)$ for 
    the sake of contradiction. Then $x \in A \cap B$ and $x \in A \backslash B$. 
    Since $x \in A \cap B$, it follows that $x \in A$ and $x \in B$. 
    Since $x \in A \backslash B$, it follows that $x \in A$ and $x \notin B$. 
    This contradicts $x \in B$. Therefore, $(A \cap B) \cap (A \backslash B) = \emptyset$.
    \item To claim $(A \cap B) \cap (B \backslash A) = \emptyset$,
    Assume there exists $x \in (A \cap B) \cap (B \backslash A)$ for 
    the sake of contradiction. Then $x \in A \cap B$ and $x \in B \backslash A$. 
    Since $x \in A \cap B$, it follows that $x \in A$ and $x \in B$. 
    Since $x \in B \backslash A$, it follows that $x \in B$ and $x \notin A$. 
    This contradicts $x \in A$. Therefore, $(A \cap B) \cap (B \backslash A) = \emptyset$.
    \item To claim $(A \backslash B) \cap (B \backslash A) = \emptyset$,
    Assume there exists $x \in (A \backslash B) \cap (B \backslash A)$ for 
    the sake of contradiction. Then $x \in A \backslash B$ and $x \in B \backslash A$. 
    Since $x \in A \backslash B$, it follows that $x \in A$ and $x \notin B$. 
    Since $x \in B \backslash A$, it follows that $x \in B$ and $x \notin A$. 
    This contradicts both $x \notin B$ and $x \notin A$ at the same time. 
    Therefore, $(A \backslash B) \cap (B \backslash A) = \emptyset$.
    \end{enumerate}
    Therefore, $(A \cap B)$, $(A \backslash B)$, and $(B \backslash A)$ are all pairwise disjoint by definition.
  \end{proof}
  \end{enumerate}
  




  \newpage
  \item Let $A, B, C$, and $D$ be subsets of a universal set $\mathcal{U}$. 
  Prove or disprove the following propositions.\\
  \textit{Likewise, Cartesian product of sets (above) and the power set of a set (below) will be very important in the next part of the course.}
  \begin{enumerate}
  \item $(A \times B) \cap(C \times D) \subseteq(A \cap C) \times(B \cap D)$.
  \begin{proof}
    Let $(x,y) \in (A \times B) \cap(C \times D)$ be arbitrary. Then $(x,y) \in A \times B$ and $(x,y) \in C \times D$. Since $(x,y) \in A \times B$, it follows that $x \in A$ and $y \in B$. Since $(x,y) \in C \times D$, it follows that $x \in C$ and $y \in D$. Therefore, $x \in A \cap C$ and $y \in B \cap D$, so $(x,y) \in (A \cap C) \times(B \cap D)$. Hence, $(A \times B) \cap(C \times D) \subseteq(A \cap C) \times(B \cap D)$.
  \end{proof}

  \item $(A \cap C) \times(B \cap D) \subseteq(A \times B) \cap(C \times D)$.
  \begin{proof}
    Let $(x,y) \in (A \cap C) \times(B \cap D)$ be arbitrary. Then $x \in A \cap C$ and $y \in B \cap D$. Since $x \in A \cap C$, it follows that $x \in A$ and $x \in C$. Since $y \in B \cap D$, it follows that $y \in B$ and $y \in D$. Therefore, $(x,y) \in A \times B$ and $(x,y) \in C \times D$, so $(x,y) \in (A \times B) \cap(C \times D)$. Hence, $(A \cap C) \times(B \cap D) \subseteq(A \times B) \cap(C \times D)$.
  \end{proof}

  \item $(A \times B) \cup(C \times D) \subseteq(A \cup C) \times(B \cup D)$.
  \begin{proof}[Disproof]
    Let $A=\{1\}, B=\{2\}, C=\{3\}, D=\{4\}$ be subsets of a universal set $\mathcal{U}$ for the sake of counterexample. Then $(A \times B) \cup(C \times D)=\{(1,2)\} \cup \{(3,4)\}=\{(1,2),(3,4)\}$ and $(A \cup C) \times(B \cup D)=\{1,3\} \times \{2,4\}=\{(1,2),(1,4),(3,2),(3,4)\}$. Since $(1,4) \in (A \cup C) \times(B \cup D)$ but $(1,4) \notin (A \times B) \cup(C \times D)$, it follows that $(A \times B) \cup(C \times D) \not\subseteq (A \cup C) \times(B \cup D)$. Therefore, the statement is \textbf{false}.
  \end{proof}

  \item $(A \cup C) \times(B \cup D) \subseteq(A \times B) \cup(C \times D)$.
  \begin{proof}[Disproof]
    Let $A=\{1\}, B=\{2\}, C=\{3\}, D=\{4\}$ be subsets of a universal set $\mathcal{U}$ for the sake of counterexample. Then $(A \cup C) \times(B \cup D)=\{1,3\} \times \{2,4\}=\{(1,2),(1,4),(3,2),(3,4)\}$ and $(A \times B) \cup(C \times D)=\{(1,2)\} \cup \{(3,4)\}=\{(1,2),(3,4)\}$. Since $(1,4) \in (A \cup C) \times(B \cup D)$ but $(1,4) \notin (A \times B) \cup(C \times D)$, it follows that $(A \cup C) \times(B \cup D) \not\subseteq (A \times B) \cup(C \times D)$. Therefore, the statement is \textbf{false}.
  \end{proof}
  \end{enumerate}





  \newpage
  \item Let $A, B$ be subsets of a universal set $\mathcal{U}$. 
  Prove or disprove the following propositions.
  \begin{enumerate}
  \item $\mathcal{P}(A \cap B) \subseteq \mathcal{P}(A) \cap \mathcal{P}(B)$.
  \begin{proof}
    Let $X \in \mathcal{P}(A \cap B)$ be arbitrary. 
    Then $X \subseteq A \cap B$ by definition of power set. 
    Since $X \subseteq A \cap B$, it follows that 
    $X \subseteq A$ and $X \subseteq B$. Therefore, $X \in \mathcal{P}(A)$ 
    and $X \in \mathcal{P}(B)$, so $X \in \mathcal{P}(A) \cap \mathcal{P}(B)$. 
    Hence, $\mathcal{P}(A \cap B) \subseteq \mathcal{P}(A) \cap \mathcal{P}(B)$.
  \end{proof}
  \item $\mathcal{P}(A) \cap \mathcal{P}(B) \subseteq \mathcal{P}(A \cap B)$.
  \begin{proof}[Disproof]
    Let $A=\{1,2\}, B=\{2,3\}$ be subsets of a universal set $\mathcal{U}$ 
    for the sake of counterexample. Then 
    $\mathcal{P}(A)=\{\emptyset,\{1\},\{2\},\{1,2\}\}$ 
    and $\mathcal{P}(B)=\{\emptyset,\{2\},\{3\},\{2,3\}\}$, 
    so $\mathcal{P}(A) \cap \mathcal{P}(B)=\{\emptyset,\{2\}\}$. 
    However, $\mathcal{P}(A \cap B)=\mathcal{P}(\{2\})=\{\emptyset,\{2\}\}$. 
    Since $\{1\} \in \mathcal{P}(A) \cap \mathcal{P}(B)$ 
    but $\{1\} \notin \mathcal{P}(A \cap B)$, 
    it follows that $\mathcal{P}(A) \cap \mathcal{P}(B) \not\subseteq \mathcal{P}(A \cap B)$. 
    Therefore, the statement is \textbf{false}.
  \end{proof}

  \item $\mathcal{P}(A \cup B) \subseteq \mathcal{P}(A) \cup \mathcal{P}(B)$.
  \begin{proof}[Disproof]
    Let $A=\{1\}, B=\{2\}$ be subsets of a universal set $\mathcal{U}$ 
    for the sake of counterexample. 
    Then $\mathcal{P}(A \cup B)=\mathcal{P}(\{1,2\})=\{\emptyset,\{1\},\{2\},\{1,2\}\}$ 
    and $\mathcal{P}(A) \cup \mathcal{P}(B)=\{\emptyset,\{1\}\} \cup \{\emptyset,\{2\}\}=\{\emptyset,\{1\},\{2\}\}$. Since $\{1,2\} \in \mathcal{P}(A \cup B)$ but $\{1,2\} \notin \mathcal{P}(A) \cup \mathcal{P}(B)$, it follows that $\mathcal{P}(A \cup B) \not\subseteq \mathcal{P}(A) \cup \mathcal{P}(B)$. 
    Therefore, the statement is \textbf{false}.
  \end{proof}

  \item $\mathcal{P}(A) \cup \mathcal{P}(B) \subseteq \mathcal{P}(A \cup B)$.
  \begin{proof}
    Let $X \in \mathcal{P}(A) \cup \mathcal{P}(B)$ be arbitrary. 
    Then $X \in \mathcal{P}(A)$ or $X \in \mathcal{P}(B)$ by definition of union. 
    If $X \in \mathcal{P}(A)$, then $X \subseteq A \subseteq A \cup B$, 
    so $X \in \mathcal{P}(A \cup B)$. If $X \in \mathcal{P}(B)$, 
    then $X \subseteq B \subseteq A \cup B$, so $X \in \mathcal{P}(A \cup B)$. 
    In both cases, $X \in \mathcal{P}(A \cup B)$. 
    Hence, $\mathcal{P}(A) \cup \mathcal{P}(B) \subseteq \mathcal{P}(A \cup B)$.
  \end{proof}
  \end{enumerate}



  
\end{enumerate}






\end{document}